% GENERATED FILE. Edit canonical lessons, then run npm run generate:books.
% canonical-source-hash: fnv1a64:c4d0a8165ba7f481
% canonical-lessons: ZH-C312-jianyu, ZH-C312-fakuan, ZH-C312-dajia, ZH-C312-huogai, ZH-C312-sihuo, ZH-R312-review-308-312

\chapter{Locked Up}
\label{ch:zh-locked-up}
\hlchaptermodality{312}

\begin{chapteropening}
\textbf{By the end of this chapter:} \emph{I can talk about fines, fights and prison, and say someone got what they deserved.}

Complete the last lesson of chapter 312: I can talk about fines, fights and prison, and say someone got what they deserved.
\end{chapteropening}

\section[jiānyù]{jiānyù — prison}

\label{lesson:ZH-C312-jianyu}



\begin{quote}
Before the new one: say the Mandarin for a warning siren, then the Mandarin for a traffic police officer.
\end{quote}

\subsection*{You'll want to know: jiānyù}
\emph{jiānyù} — \textquotedblleft{}prison\textquotedblright{}.

\begin{quote}
\emph{tā zài jiānyù lǐ zhù le sān nián} — He spent three years in prison.
\end{quote}

Its characters come later; for now it is a word for the ear.

\begin{cousinweb}
\emph{jiānyù} is built from \emph{jiān} 'to watch over' and \emph{yù} 'a jail', both inherited from Middle Chinese, \emph{kaem} and \emph{ngjowk} (Guangyun, in Baxter's transcription), which glosses the first 'to oversee' and explains the shape of the second as two dogs on guard.
\end{cousinweb}

\begin{grammarlens}[title={What you've built}]
One more word for telling what happened.
\end{grammarlens}

\subsection*{Your turn}
\begin{itemize}
\raggedright
  \item \emph{Say it:} \emph{jiānyù}
  \item \emph{Say it:} \emph{jiānyù}, once more
  \item \emph{Say it:} \emph{tā zài jiānyù lǐ zhù le sān nián} — He spent three years in prison.
  \item \emph{Recall:} say the Mandarin for a warning siren, then the Mandarin for a traffic police officer, then say \emph{jiānyù} again
\end{itemize}

\begin{tcolorbox}[breakable,colback=teal!4,colframe=teal!35!black,title={Before you move on}]
What does jiānyù mean? (\textquotedblleft{}Prison\textquotedblright{}.) Say it once more.
\end{tcolorbox}
\section[fákuǎn]{fákuǎn — to impose a fine; a penalty fine}

\label{lesson:ZH-C312-fakuan}



\begin{quote}
Before the new one: say the Mandarin for a traffic police officer, then the Mandarin for a prison.
\end{quote}

\subsection*{You'll want to know: fákuǎn}
\emph{fákuǎn} — \textquotedblleft{}to impose a fine; a penalty fine\textquotedblright{}.

\begin{quote}
\emph{kāichē tài kuài yào fákuǎn} — Driving too fast gets you a fine.
\end{quote}

Its characters come later; for now it is a word for the ear.

\begin{cousinweb}
\emph{fákuǎn} is built from \emph{fá} 'to punish' and \emph{kuǎn} 'a sum', both inherited from Middle Chinese, \emph{bjot} and \emph{khwanX} (Guangyun, in Baxter's transcription), which glosses the first 'a punishment' and the second 'sincere'; its sense 'a sum of money' is not given there.
\end{cousinweb}

\begin{grammarlens}[title={What you've built}]
One more word for telling what happened.
\end{grammarlens}

\subsection*{Your turn}
\begin{itemize}
\raggedright
  \item \emph{Say it:} \emph{fákuǎn}
  \item \emph{Say it:} \emph{fákuǎn}, once more
  \item \emph{Say it:} \emph{kāichē tài kuài yào fákuǎn} — Driving too fast gets you a fine.
  \item \emph{Recall:} say the Mandarin for a traffic police officer, then the Mandarin for a prison, then say \emph{fákuǎn} again
\end{itemize}

\begin{tcolorbox}[breakable,colback=teal!4,colframe=teal!35!black,title={Before you move on}]
What does fákuǎn mean? (\textquotedblleft{}To impose a fine; a penalty fine\textquotedblright{}.) Say it once more.
\end{tcolorbox}
\section[dǎjià]{dǎjià — to fight; to get into a fistfight}

\label{lesson:ZH-C312-dajia}



\begin{quote}
Before the new one: say the Mandarin for a prison, then the Mandarin for to impose a fine.
\end{quote}

\subsection*{You'll want to know: dǎjià}
\emph{dǎjià} — \textquotedblleft{}to fight; to get into a fistfight\textquotedblright{}.

\begin{quote}
\emph{liǎng ge rén zhèngzài ménkǒu dǎjià} — Two people were fighting at the door.
\end{quote}

Its characters come later; for now it is a word for the ear.

\begin{cousinweb}
\emph{dǎjià} is built from \zh{打} 'to hit' and \emph{jià} 'a frame', both inherited from Middle Chinese, \emph{taengX} and \emph{kaeH} (Guangyun, in Baxter's transcription), which glosses the first 'to strike' and the second 'a roof frame'; how the pair came to mean a fight is not settled, and dǎ is an irregular outcome of \emph{taengX}.
\end{cousinweb}

\begin{grammarlens}[title={What you've built}]
One more word for telling what happened.
\end{grammarlens}

\subsection*{Your turn}
\begin{itemize}
\raggedright
  \item \emph{Say it:} \emph{dǎjià}
  \item \emph{Say it:} \emph{dǎjià}, once more
  \item \emph{Say it:} \emph{liǎng ge rén zhèngzài ménkǒu dǎjià} — Two people were fighting at the door.
  \item \emph{Recall:} say the Mandarin for a prison, then the Mandarin for to impose a fine, then say \emph{dǎjià} again
\end{itemize}

\begin{tcolorbox}[breakable,colback=teal!4,colframe=teal!35!black,title={Before you move on}]
What does dǎjià mean? (\textquotedblleft{}To fight; to get into a fistfight\textquotedblright{}.) Say it once more.
\end{tcolorbox}
\section[\texorpdfstring{\zh{活该} (huógāi)}{活该 (huógāi)}]{\zh{活该} (huógāi) — it serves them right; serves you right}

\label{lesson:ZH-C312-huogai}



\begin{quote}
Before the new one: say the Mandarin for to impose a fine, then the Mandarin for to get into a fistfight.
\end{quote}

\subsection*{You'll want to know: \zh{活该}}
\textbf{\zh{活该}} — \emph{huógāi} — \textquotedblleft{}it serves them right; serves you right\textquotedblright{}.

\begin{quote}
\textbf{\zh{活该他经常偷东西}} \emph{huógāi, tā jīngcháng tōu dōngxi} — Serves him right, he steals all the time.
\end{quote}

\begin{cousinweb}
\textbf{\zh{活该}} (\emph{huógāi}) is built from \zh{活} 'alive' and \zh{该} 'ought', both inherited from Middle Chinese, \emph{hwat} and \emph{koj} (Guangyun, in Baxter's transcription), which glosses the second 'complete' and 'all'; how the pair came to mean 'it serves you right' is not settled.
\end{cousinweb}

\begin{grammarlens}[title={What you've built}]
One more word for telling what happened.
\end{grammarlens}

\subsection*{Your turn}
\begin{itemize}
\raggedright
  \item \emph{Say it:} \emph{huógāi}
  \item \emph{Say it:} \emph{huógāi}, once more
  \item \emph{Say it:} \emph{huógāi, tā jīngcháng tōu dōngxi} — Serves him right, he steals all the time.
  \item \emph{Recall:} say the Mandarin for to impose a fine, then the Mandarin for to get into a fistfight, then say \emph{huógāi} again
\end{itemize}

\begin{tcolorbox}[breakable,colback=teal!4,colframe=teal!35!black,title={Before you move on}]
What does \zh{活该} mean? (\textquotedblleft{}It serves them right; serves you right\textquotedblright{}.) Say it once more.
\end{tcolorbox}
\section[\texorpdfstring{\zh{死活} (sǐhuó)}{死活 (sǐhuó)}]{\zh{死活} (sǐhuó) — no matter what; whatever anyone says}

\label{lesson:ZH-C312-sihuo}



\begin{quote}
Before the new one: say the Mandarin for to get into a fistfight, then the Mandarin for it serves them right.
\end{quote}

\subsection*{You'll want to know: \zh{死活}}
\textbf{\zh{死活}} — \emph{sǐhuó} — \textquotedblleft{}no matter what; whatever anyone says\textquotedblright{}.

\begin{quote}
\textbf{\zh{他死活不去报警}} \emph{tā sǐhuó bú qù bàojǐng} — Whatever anyone said, he wouldn't go to the police.
\end{quote}

\begin{cousinweb}
\textbf{\zh{死活}} (\emph{sǐhuó}) is built from \zh{死} 'to die' and \zh{活} 'to live', both inherited from Middle Chinese, \emph{sijX} and \emph{hwat} (Guangyun, in Baxter's transcription): dead or alive, whatever happens.
\end{cousinweb}

\begin{grammarlens}[title={What you've built}]
One more word for telling what happened.
\end{grammarlens}

\subsection*{Your turn}
\begin{itemize}
\raggedright
  \item \emph{Say it:} \emph{sǐhuó}
  \item \emph{Say it:} \emph{sǐhuó}, once more
  \item \emph{Say it:} \emph{tā sǐhuó bú qù bàojǐng} — Whatever anyone said, he wouldn't go to the police.
  \item \emph{Recall:} say the Mandarin for to get into a fistfight, then the Mandarin for it serves them right, then say \emph{sǐhuó} again
\end{itemize}

\begin{tcolorbox}[breakable,colback=teal!4,colframe=teal!35!black,title={Before you move on}]
What does \zh{死活} mean? (\textquotedblleft{}No matter what; whatever anyone says\textquotedblright{}.) Say it once more.
\end{tcolorbox}
\section[(dialogue)]{Looking back: from Ever Since to Locked Up}

\label{lesson:ZH-R312-review-308-312}



\begin{quote}
Say the Mandarin for to go missing and ever since.
\end{quote}

\subsection*{Your turn}
\begin{itemize}
\raggedright
  \item \emph{Say it:} to go missing, ever since, to come from a place, to commit a crime, the local police station
  \item \emph{Say it:} to rob by force, a thief, to con someone, a con man, on purpose
  \item \emph{Say it:} to suspect someone, a witness to a crime, proof of what happened, to tell a lie, to own up to something
  \item \emph{Say it:} to investigate a case, a police car, to run after someone, a warning siren, a traffic police officer
  \item \emph{Say it:} a prison, to impose a fine, to get into a fistfight, it serves them right, whatever anyone says
  \item \emph{Read it:} \zh{自} \zh{偷} \zh{警} — each one aloud, then a word that holds it
\end{itemize}

\begin{tcolorbox}[breakable,colback=teal!4,colframe=teal!35!black,title={Before you move on}]
To go missing? (\emph{shīzōng}.) Ever since? (\textbf{\zh{自从}}, \emph{zìcóng}.) To come from a place? (\textbf{\zh{来自}}, \emph{láizì}.) To commit a crime? (\emph{fànzuì}.) The local police station? (\emph{pàichūsuǒ}.) To rob by force? (\emph{qiǎngjié}.) A thief? (\textbf{\zh{小偷}}, \emph{xiǎotōu}.) To con someone? (\emph{piàn}.) A con man? (\emph{piànzi}.) On purpose? (\textbf{\zh{故意}}, \emph{gùyì}.) To suspect someone? (\emph{huáiyí}.) A witness to a crime? (\emph{zhèngrén}.) Proof of what happened? (\emph{zhèngjù}.) To tell a lie? (\emph{shuōhuǎng}.) To own up to something? (\emph{chéngrèn}.) To investigate a case? (\emph{diàochá}.) A police car? (\textbf{\zh{警车}}, \emph{jǐngchē}.) To run after someone? (\emph{zhuī}.) A warning siren? (\textbf{\zh{警报}}, \emph{jǐngbào}.) A traffic police officer? (\emph{jiāojǐng}.) A prison? (\emph{jiānyù}.) To impose a fine? (\emph{fákuǎn}.) To get into a fistfight? (\emph{dǎjià}.) It serves them right? (\textbf{\zh{活该}}, \emph{huógāi}.) Whatever anyone says? (\textbf{\zh{死活}}, \emph{sǐhuó}.)
\end{tcolorbox}
